\chapter{1000 Câu Cà Khịa Học Đường}
\markboth{1000 Câu Cà Khịa Học Đường}{1000 Câu Cà Khịa Học Đường}

\begin{truyen}{Câu 1–100 trong bộ 1000 câu}{Entries 1–100 of the 1000-entry collection}
\chuhoa{H}{ọc đường không chỉ là nơi nhồi nhét kiến thức, mà còn là đấu trường của những sự thật không ai chịu nói thẳng. Chương này tập hợp những câu cà khịa --- không phải để chế giễu, mà để mỗi người đọc xong tự hỏi xem mình có trong đó không. Nếu có thì chúc mừng: em vừa được soi miễn phí.}
\textit{(School is not merely a place for stuffing knowledge; it is an arena of truths no one wants to say plainly. This chapter collects roasting lines --- not to mock, but so each reader wonders whether they appear in it. If you do: congratulations, you just received a free mirror check.)}
\end{truyen}

%------------------------------------------------------------
\section{Cà Khịa Bằng Số Liệu và Logic Không Thể Phủ Nhận (Câu 1–25)}
%------------------------------------------------------------
\begin{enumerate}[leftmargin=*]

\songnguitem{Nghiên cứu cho thấy 100\% học sinh không làm bài tập đều có lý do rất chính đáng. Nghiên cứu đó do chính học sinh đó tự làm ra.}{Research shows 100\% of students who do not do homework have very valid reasons. The research was conducted by those same students.}

\songnguitem{Tỷ lệ mở sách đêm trước khi thi tỉ lệ nghịch với điểm số nhận được sáng hôm sau, nhưng học sinh vẫn thử điều này mỗi năm.}{The rate of opening textbooks the night before an exam is inversely proportional to the scores received the next morning, yet students keep trying this approach every year.}

\songnguitem{Có một thống kê không chính thức: lý do phổ biến nhất để không làm bài là ``quên'', đứng thứ hai là ``chó ăn mất'', thứ ba là ``không biết là có bài''.}{There is an unofficial statistic: the most common excuse for not doing homework is `forgot'; second is `the dog ate it'; third is `didn't know there was any'.}

\songnguitem{Thời gian một học sinh dành để tìm cách chép bài của bạn thường dài hơn thời gian cần để tự làm bài đó.}{The time a student spends trying to copy someone else's homework is usually longer than the time it would take to simply do it themselves.}

\songnguitem{Bài văn được viết trong mười lăm phút cuối giờ trông khá giống bài được viết cả tuần, miễn là không đọc kỹ.}{An essay written in the last fifteen minutes looks very similar to one written over a whole week --- as long as no one reads it carefully.}

\songnguitem{Tỷ lệ học sinh nói ``em hiểu rồi'' nhưng thực ra chưa hiểu là con số mà không một bảng thống kê nào dám công bố.}{The percentage of students who say `I understand now' while actually not understanding is a figure no statistical chart dares to publish.}

\songnguitem{Khoảng cách giữa việc muốn học giỏi và việc thực sự ngồi xuống học có đơn vị đo là ``tháng''.}{The distance between wanting to study well and actually sitting down to study is measured in months.}

\songnguitem{Ghế trong lớp đang cạnh tranh khốc liệt với giường ngủ để giành vị trí ``nơi học sinh ngủ nhiều nhất''.}{The classroom chair is in fierce competition with the bed for the title of `place where students sleep the most'.}

\songnguitem{Câu hỏi ``học để làm gì?'' đã được hỏi bởi mọi thế hệ học sinh, nhưng chỉ được trả lời đúng bởi những người đã tốt nghiệp và muốn quay lại.}{The question `What is studying for?' has been asked by every generation of students, but only answered accurately by those who have graduated and wish they could return.}

\songnguitem{Có một quy luật vật lý không ghi trong sách giáo khoa: cuốn sách càng nặng thì học sinh càng không mang đến lớp.}{There is an unwritten law of physics: the heavier the textbook, the less likely a student is to bring it to class.}

\songnguitem{Điện thoại của học sinh sạc pin nhanh hơn bộ não được nạp kỹ năng trong lớp học.}{A student's phone charges faster than the student's brain absorbs skills during class.}

\songnguitem{Đề thi không bao giờ dễ như khi học sinh tưởng đêm trước, và không bao giờ khó như học sinh than thở ngay sau khi thi xong.}{Exams are never as easy as students imagine the night before, and never as hard as they complain about immediately after finishing.}

\songnguitem{Ba điều chắc chắn trong cuộc đời: cái chết, thuế, và học sinh sẽ tìm ra cách ngồi gần bạn giỏi nhất trong giờ kiểm tra.}{Three things are certain in life: death, taxes, and students finding a way to sit next to the best student during an exam.}

\songnguitem{Trí nhớ học sinh hoạt động theo nguyên tắc: nhớ hết mọi thứ không liên quan đến bài thi cho đến đúng lúc ngồi vào phòng thi thì quên sạch những gì cần nhớ.}{A student's memory operates on a principle: remember everything unrelated to the exam until the exact moment of sitting down in the exam room, then forget everything that matters.}

\songnguitem{Tỷ lệ học sinh lắng nghe khi thầy giải thích bài khó tỉ lệ nghịch với độ khó của bài đó.}{The proportion of students listening when a teacher explains a difficult topic is inversely proportional to the difficulty of that topic.}

\songnguitem{Giờ Văn có một hiện tượng kỳ lạ: học sinh không đọc sách nhưng có thể phân tích nhân vật rất sâu sắc dựa trên tóm tắt hai đoạn.}{Literature class has a strange phenomenon: students who haven't read the book can analyze characters at great depth based on a two-paragraph summary.}

\songnguitem{Một học sinh giải thích bài toán sai qua ba giai đoạn: tự tin, bướng bỉnh, và cuối cùng là ``à em nhớ ra rồi''.}{A student explaining an incorrect solution goes through three stages: confidence, stubbornness, and finally `oh wait, I remember now'.}

\songnguitem{Xin phép đi vệ sinh trong giờ thi là một chiến thuật cổ điển có hiệu quả giảm dần theo từng năm nhưng vẫn tiếp tục được sử dụng.}{Asking to go to the restroom during an exam is a classic strategy with diminishing effectiveness each year, yet it continues to be deployed.}

\songnguitem{Cứ mỗi khi thầy cô nói ``đây không ra trong thi đâu'' là học sinh chú ý hơn bất kỳ lúc nào khác trong học kỳ.}{Whenever a teacher says `this won't be on the test', students pay more attention than at any other point in the semester.}

\songnguitem{Không ai quan tâm đến số trang sách giáo khoa cho đến khi thầy cô nói ``học từ trang này đến trang kia''.}{No one pays attention to textbook page numbers until the teacher announces `study from this page to that one'.}

\songnguitem{Bài thuyết trình được chuẩn bị tốt nhất là bài được chuẩn bị tối hôm trước khi trình bày, điều mà không ai thừa nhận nhưng tất cả đều biết.}{The best-prepared presentations are prepared the night before they are due --- something no one admits but everyone knows.}

\songnguitem{Học sinh hay nói ``em không có năng khiếu'', nhưng thực ra chỉ cần thêm mười hai tiếng với cuốn sách là năng khiếu đó sẽ xuất hiện.}{Students often say `I have no talent,' when in reality twelve more hours with a textbook would make that talent appear.}

\songnguitem{Điểm số tốt và điện thoại thông minh nằm trong cùng một phương trình, và quan hệ của chúng thường là tỉ lệ nghịch.}{Good grades and smartphones exist in the same equation, and their relationship is usually inversely proportional.}

\songnguitem{Quy tắc ngầm của học đường: người nộp bài đầu tiên không phải lúc nào cũng giỏi nhất, nhưng người nộp bài cuối thường biết lý do vì sao.}{The unspoken rule of school: the student who submits first isn't always the best, but the one who submits last usually knows exactly why.}

\songnguitem{Học sinh nào cũng tin mình học tốt hơn kết quả thể hiện. Thống kê cho thấy không phải vậy. Học sinh vẫn tin.}{Every student believes they performed better than their results show. Statistics say otherwise. Students continue believing.}

\end{enumerate}

%------------------------------------------------------------
\section{Giữa Mơ và Thực Tế Là Một Vực Thẳm Rất Sâu (Câu 26–50)}
%------------------------------------------------------------
\begin{enumerate}[leftmargin=*,resume]

\songnguitem{Giấc mơ làm ca sĩ rất to, nhưng bài nhạc lý điểm ba thì cần được nhắc đến.}{The dream of becoming a singer is grand --- but the music theory score of three deserves a mention.}

\songnguitem{Ai cũng muốn làm bác sĩ nhưng không ai muốn học Sinh, điều này tạo ra một sự mâu thuẫn rất thú vị để suy ngẫm.}{Everyone wants to become a doctor, but no one wants to study Biology --- this creates a very interesting contradiction to reflect on.}

\songnguitem{Ước mơ làm triệu phú không có gì sai, nhưng ước mơ đó cần một kế hoạch không phải là ``trúng số''.}{Dreaming of becoming a millionaire is not wrong, but that dream needs a plan that does not involve `winning the lottery'.}

\songnguitem{Có những học sinh viết trong bài văn ``em muốn trở thành nhà khoa học'' trong khi môn Vật lý vẫn đang ở mức ba điểm.}{Some students write in essays `I want to become a scientist' while their Physics score is still sitting at three points.}

\songnguitem{Không ai nói với học sinh rằng khoảng cách giữa ``tôi sẽ học sau'' và ``tôi đã không học'' là bằng không.}{No one tells students that the distance between `I'll study later' and `I didn't study' is exactly zero.}

\songnguitem{Giấc mơ lớn không có gì sai, vấn đề là nhiều người cư xử như giấc mơ đó sẽ tự hiện thực mà không cần họ làm gì.}{Having big dreams is not wrong; the problem is that many people act as though those dreams will come true without any effort from them.}

\songnguitem{Học sinh hay nói ``em không phù hợp với trường này'' khi thực ra vấn đề là trường yêu cầu học bài còn học sinh thì không muốn.}{Students often say `I'm not suited for this school' when the real issue is that the school requires studying and the student simply does not want to.}

\songnguitem{Ước mơ làm YouTuber rất phổ biến. Ít ai nhắc đến tỷ lệ người thực sự thành công so với số người thử.}{The dream of becoming a YouTuber is very popular. Few mention the ratio of those who actually succeed to those who try.}

\songnguitem{Kế hoạch thi đại học tốt nhất mà nhiều học sinh đang theo đuổi là: ``tính sau''.}{The best university entrance plan many students are currently pursuing is: `figure it out later'.}

\songnguitem{Học sinh nói ``em học theo cách của em'' đôi khi là sự thật, đôi khi là cách nói khác của ``em không học''.}{`I study in my own way' is sometimes true and sometimes simply another way of saying `I don't study'.}

\songnguitem{Thế hệ trước nói ``chịu khó học đi'' và thế hệ này trả lời ``học không bằng có tài năng''. Không ai trong số đó làm cả hai.}{The previous generation said `work hard at your studies' and this generation replied `hard work doesn't beat talent'. Neither group managed to do both.}

\songnguitem{Tương lai tươi sáng mà học sinh thấy trước mắt thường rất đẹp và rất mơ hồ --- hai điều này hay đi cùng nhau.}{The bright future students see ahead is usually very beautiful and very vague --- two things that tend to go together.}

\songnguitem{Nhiều người nói ``học không bằng được việc'', điều đó đúng với một số ít trường hợp và được trích dẫn bởi 100\% học sinh lười.}{Many people say `success doesn't require a degree,' which is true for a small minority and quoted by 100\% of students who don't want to study.}

\songnguitem{Sự tự tin của học sinh vào khả năng vượt qua kỳ thi mà không cần chuẩn bị là một hiện tượng tâm lý cần được nghiên cứu kỹ.}{Student confidence in passing exams without preparation is a psychological phenomenon deserving of serious study.}

\songnguitem{Có một cái bẫy rất phổ biến trong học đường: tin rằng điểm số thấp không phản ánh năng lực thật, trong khi không làm gì để chứng minh điều đó.}{There is a very common trap in school: believing that low scores don't reflect your true ability, while doing nothing to prove it.}

\songnguitem{Học sinh giỏi viết mục tiêu trong vở. Học sinh chưa giỏi viết mục tiêu lên mạng xã hội và không làm gì thêm.}{Hardworking students write their goals in notebooks. Struggling students post their goals on social media and do nothing further.}

\songnguitem{Nhiều bạn trẻ muốn được tôn trọng mà không muốn làm điều gì để xứng đáng với sự tôn trọng đó.}{Many young people want to be respected without wanting to do anything to deserve it.}

\songnguitem{``Áp lực'' là từ nhiều học sinh dùng cho cả tình trạng thực sự khó khăn lẫn tình trạng bị yêu cầu làm bài tập về nhà.}{`Pressure' is a word many students use both for genuinely difficult circumstances and for simply being asked to do their homework.}

\songnguitem{Học sinh nào cũng nghĩ mình sẽ chăm hơn ở học kỳ sau. Mỗi học kỳ.}{Every student thinks they will work harder next semester. Every semester.}

\songnguitem{Câu ``em sẽ cố gắng'' trong học đường thường có vòng đời hai tuần, sau đó trở về trạng thái ban đầu.}{The sentence `I will try harder' in a school context usually has a two-week lifespan before returning to its original state.}

\songnguitem{Kỳ nghỉ hè là lúc nhiều học sinh lên kế hoạch học bù, và cũng là lúc họ không thực hiện điều đó.}{Summer break is when many students plan to catch up on studying, and also when they don't do it.}

\songnguitem{Học sinh hay hỏi ``học cái này để làm gì?'' Người trả lời được câu đó thường đã học xong và đang đi làm.}{Students often ask `What is this learning for?' The people who can answer that question have usually already graduated and are working.}

\songnguitem{Việc chuẩn bị cho tương lai thường ít thú vị hơn việc nghĩ về tương lai --- và đây là vấn đề cốt lõi.}{Preparing for the future is usually less enjoyable than thinking about the future --- and this is the core problem.}

\songnguitem{Không có ai thành công nhờ ước mơ mà không có kỷ luật. Câu này đúng nhưng ít ai thích nghe.}{No one succeeds on dreams without discipline. This is true, but rarely welcomed.}

\songnguitem{Học sinh nào cũng muốn thay đổi. Ít người muốn chịu đựng quá trình thay đổi đó.}{Every student wants change. Few are willing to endure the process of changing.}

\end{enumerate}

%------------------------------------------------------------
\section{An Ủi Theo Cách Không Ai Muốn Được An Ủi (Câu 51–75)}
%------------------------------------------------------------
\begin{enumerate}[leftmargin=*,resume]

\songnguitem{Đừng buồn vì bị điểm thấp. Ít nhất em biết mình đang đứng ở đâu, và đó là thông tin rất có giá trị.}{Don't be sad about a low score. At least you know where you stand, and that is very valuable information.}

\songnguitem{Thất bại là bài học. Vậy nên với số lần thất bại của em, em đang học được rất nhiều.}{Failure is a lesson --- so with your number of failures, you must be learning a great deal.}

\songnguitem{Em không phải người duy nhất chưa vượt qua bài kiểm tra này. Em chỉ là người duy nhất chưa vượt qua nó lần thứ tư.}{You are not the only person who hasn't passed this test --- you are just the only one who hasn't passed it for the fourth time.}

\songnguitem{Điểm thấp không nói lên tất cả. Nó chỉ nói lên phần lớn.}{A low score does not say everything. Just most of it.}

\songnguitem{Đừng nản, đường còn dài và em còn nhiều cơ hội để làm lại hầu hết những gì em đã làm sai.}{Don't give up --- the road is long and you have many opportunities to redo most of what you have done wrong.}

\songnguitem{Em còn trẻ, còn có thời gian để sửa. Nhưng thời gian đó không nhiều như em đang nghĩ.}{You are young and still have time to improve. But not as much time as you think.}

\songnguitem{Ai cũng có những bài thi không như mong muốn. Điểm của em chỉ là một ví dụ đặc biệt ấn tượng.}{Everyone has exams that don't go as hoped. Your grade is just a particularly memorable example.}

\songnguitem{Em không phải người kém nhất lớp. Nhưng thứ hạng đó đang bị thách thức mỗi ngày.}{You are not the weakest student in class. But that position is being challenged every single day.}

\songnguitem{Điểm số chỉ là một con số, nhưng con số đó đang cố nói với em điều gì đó rất quan trọng.}{A score is just a number --- but that number is trying to tell you something very important.}

\songnguitem{Đừng lo khi không ai tin em có thể làm được. Đó là cơ hội để em tự chứng minh --- hoặc không.}{Don't worry when no one believes you can do it. That is an opportunity to prove them wrong. Or not.}

\songnguitem{Kết quả của em hôm nay phản ánh nỗ lực của em tháng trước. Và điều đó giải thích rất nhiều.}{Today's results reflect last month's effort. And that explains a great deal.}

\songnguitem{Thất bại không định nghĩa em, nhưng cách em phản ứng với thất bại thì có --- và hiện tại cách em đang phản ứng là ngủ thêm.}{Failure doesn't define you, but how you respond to failure does --- and right now your response is sleeping more.}

\songnguitem{Không có gì sai khi điểm thấp. Cái có vấn đề là điểm thấp và không thay đổi gì cả.}{There is nothing wrong with a low score by itself. The problem is a low score and absolutely no changes.}

\songnguitem{Em sẽ nhìn lại kết quả này sau vài năm và thấy nó buồn cười. Còn bây giờ chỉ cần thấy nó cần sửa là đủ.}{You will look back at these results in a few years and find them funny. For now, just finding them fixable is enough.}

\songnguitem{Không phải em tệ, chỉ là em chưa tốt --- và khoảng cách giữa hai cái đó phụ thuộc vào việc em có chịu thu hẹp nó không.}{You are not bad --- you are just not good yet. The distance between those two things depends on whether you choose to close it.}

\songnguitem{Đừng so sánh mình với bạn giỏi hơn theo kiểu tức tối. Hãy so sánh theo kiểu hỏi bạn đó làm khác mình cái gì.}{Don't compare yourself to stronger students out of frustration. Compare yourself by asking what they do differently from you.}

\songnguitem{Thầy vẫn tin vào em, nhưng niềm tin của thầy đang cần thêm bằng chứng để tiếp tục tồn tại.}{I still believe in you --- but my belief is in need of more evidence to survive.}

\songnguitem{Em làm được, thầy chắc vậy. Chỉ là thầy không chắc khi nào.}{You can do it --- I am sure of that. I am just not sure of when.}

\songnguitem{Điều quan trọng không phải em đang ở đâu mà là em đang đi về đâu. Nhưng nhìn hướng em đang đi thì cũng lo một chút.}{What matters is not where you are but where you are going --- though the direction you are currently heading does cause some concern.}

\songnguitem{Hãy coi điểm thấp là bản đồ. Nó đang chỉ ra nơi em cần đến nhưng chưa đến được.}{Think of low scores as a map pointing to where you need to go but haven't yet reached.}

\songnguitem{Cuộc đời dài lắm, em còn nhiều cơ hội --- nhưng mỗi cơ hội bị bỏ qua thì số cơ hội còn lại cũng giảm theo.}{Life is long and you have many chances --- but each wasted opportunity also reduces the number remaining.}

\songnguitem{Em không nên nản, nhưng cũng không nên thoải mái. Cái thầy muốn em có là thứ nằm giữa hai cái đó: lo đúng chỗ.}{You shouldn't despair, but you also shouldn't be comfortable. What I want you to have is something between the two: worry in the right places.}

\songnguitem{Thầy hi vọng kết quả này là cú đánh thức em cần. Nếu không thì kết quả tiếp theo sẽ phải đảm nhận nhiệm vụ đó.}{I hope this result is the wake-up call you need. If not, the next result will have to take on that responsibility.}

\songnguitem{Điểm số tệ không phải là dấu chấm hết, nhưng nếu em không làm gì thì đó là dấu phẩy trước một chuỗi dài từ ``tiếp tục tệ''.}{A bad score is not the end --- but if you do nothing, it is a comma before a long string of `continued bad scores'.}

\songnguitem{Thầy chúc em mau tiến bộ, vì thầy đã chúc điều đó ba học kỳ liên tiếp và đang tự hỏi bao giờ lời chúc đó thành hiện thực.}{I wish you rapid improvement --- I have been making this wish for three consecutive semesters and am starting to wonder when it will come true.}

\end{enumerate}

%------------------------------------------------------------
\section{Quan Sát Học Sinh Trong Môi Trường Tự Nhiên (Câu 76–100)}
%------------------------------------------------------------
\begin{enumerate}[leftmargin=*,resume]

\songnguitem{Học sinh trong giờ thầy giảng bài và học sinh trong giờ ra chơi là hai sinh vật hoàn toàn khác nhau về mức độ năng lượng.}{A student during a lecture and the same student during recess are two completely different creatures in terms of energy levels.}

\songnguitem{Giờ học cuối trước kỳ nghỉ và giờ học đầu sau kỳ nghỉ là hai khung thời gian giáo viên không nên kỳ vọng quá nhiều.}{The last class before a holiday and the first class after a holiday are two time windows when teachers should not expect too much.}

\songnguitem{Một học sinh nhìn ra cửa sổ trong giờ học đang thực hiện một trong hai việc: mơ mộng về tương lai hoặc đếm xe.}{A student staring out the window during class is doing one of two things: dreaming about the future or counting passing vehicles.}

\songnguitem{Học sinh hay trả lời ``em không biết'' theo ba giọng điệu: thật sự không biết, biết nhưng không muốn nói, và biết nhưng chưa nghĩ ra.}{Students answer `I don't know' in three tones: genuinely don't know, know but don't want to say, and know but haven't thought it through yet.}

\songnguitem{Nhóm ngồi bàn cuối lớp và nhóm ngồi bàn đầu lớp thường có quan điểm rất khác nhau về mục đích của việc đến trường.}{The group at the back and the group at the front usually have very different views on the purpose of coming to school.}

\songnguitem{Học sinh giơ tay phát biểu và học sinh phát biểu vì bị gọi tên có hai nét mặt khác nhau rất rõ ràng.}{Students who raise their hands voluntarily and students who are called on have two very distinctly different facial expressions.}

\songnguitem{Giờ kiểm tra là lúc duy nhất trong năm học mà tất cả học sinh đều im lặng và tập trung. Chỉ tiếc là vào thứ họ đang tập trung.}{Exam time is the only moment in the school year when all students are simultaneously silent and focused. Unfortunately the focus is the problem.}

\songnguitem{Cuốn sách giáo khoa mới nhất và sạch nhất trong lớp thường thuộc về học sinh ít khi dùng đến nó nhất.}{The newest and cleanest textbook in class usually belongs to the student who uses it the least.}

\songnguitem{Học sinh nào cũng biết đáp án ``đúng'' khi nhìn vào bài của người khác sau khi đã nộp bài của mình.}{Every student knows the `right' answer when looking at someone else's paper after submitting their own.}

\songnguitem{Thời điểm học sinh chăm chỉ nhất trong năm thường là ngay sau khi thầy cô thông báo điểm học kỳ sắp đóng.}{The hardest-working moment in a student's year is right after the teacher announces grades are about to be finalized.}

\songnguitem{Khi thầy nói ``ai chưa hiểu giơ tay'', cả lớp đều im lặng, kể cả những người thực sự chưa hiểu.}{When a teacher says `raise your hand if you don't understand', the whole class stays silent --- including those who genuinely don't.}

\songnguitem{Sách giáo khoa Toán có nhiều bụi nhất trong các loại sách --- đây là dữ liệu không chính thức nhưng rất đáng tin cậy.}{Mathematics textbooks have the most dust of all school books --- unofficial data, but highly reliable.}

\songnguitem{Học sinh thường chú ý nhất đến câu hỏi của thầy cô ngay sau khi giật mình vì câu hỏi đó không liên quan đến bài đang học.}{Students pay closest attention to a teacher's question right after jumping in surprise because it has nothing to do with the current topic.}

\songnguitem{Điểm danh là thủ tục duy nhất mà học sinh luôn trả lời to và rõ ràng, dù cả giờ học vừa rồi không nói gì nghe được.}{Roll call is the only procedure where students always respond loudly and clearly, even if they said nothing audible for the entire preceding class.}

\songnguitem{Học sinh xin nghỉ vì lý do sức khỏe và học sinh xin nghỉ vì không muốn đi thường viết phiếu xin phép giống nhau một cách đáng kinh ngạc.}{Students absent for health reasons and those absent because they don't want to attend tend to write surprisingly similar excuse notes.}

\songnguitem{Bài văn cảm nhận về cuốn sách học sinh chưa đọc là một thể loại văn học ít được nhắc đến nhưng rất phong phú.}{An essay reflecting on a book a student has not read is a literary genre rarely discussed but extremely rich in variety.}

\songnguitem{Câu ``Thầy ơi, bao giờ thi?'' thường được hỏi đúng ngay sau khi thầy vừa thông báo lịch thi xong.}{`Teacher, when is the exam?' is most commonly asked immediately after the teacher just announced the exam schedule.}

\songnguitem{Học sinh giỏi giúp bạn hiểu bài là điều đẹp. Học sinh giỏi để bạn chép bài thì khác hẳn. Hai việc này hay bị nhầm lẫn.}{A strong student helping others understand is admirable. A strong student letting others copy is different entirely. The two are frequently confused.}

\songnguitem{Kỳ nghỉ dài nhất trong năm thường là lúc học sinh nghĩ ra lý do tại sao mình chưa tiến bộ và đổ lỗi cho kỳ nghỉ đó.}{The longest holiday of the year is usually when students think up reasons why they haven't improved and blame it on the holiday itself.}

\songnguitem{Học sinh hay nói ``em không quen học nhiều một lúc'' nhưng lại hoàn toàn quen xem phim nhiều tiếng liên tiếp không nghỉ.}{Students often say `I'm not used to studying for long stretches' while being perfectly accustomed to watching shows for hours without stopping.}

\songnguitem{Bài thuyết trình đầu năm và bài thuyết trình cuối năm của cùng một học sinh cách nhau không đáng kể nếu không có phản hồi ở giữa.}{A student's first presentation and their last presentation of the year differ very little if no feedback was given in between.}

\songnguitem{Có một quy luật ngầm: càng không học thì lý do để không học càng nhiều và càng được bản thân tin là có thật.}{There is an unwritten rule: the less you study, the more reasons you accumulate for not studying --- and the more convincingly you believe those reasons.}

\songnguitem{Giờ tự học trong lớp là thời gian mà ``tự'' thường được thực hiện đúng nghĩa, còn ``học'' thì tùy.}{Study hall is a period where `self' is usually respected as intended, while `study' remains entirely optional.}

\songnguitem{Học sinh tốt nhất và học sinh cần cố gắng nhất trong lớp có một điểm chung: cả hai đều đang ngồi trong phòng học này ngay lúc này.}{The best student and the student who needs to try hardest share one thing in common: both are sitting in this classroom right now.}

\songnguitem{Trường học là nơi kỳ lạ nhất: nơi người ta đến để học nhưng dành nhiều năng lượng nhất cho việc tìm cách không học.}{School is the strangest place: people come to learn but spend most of their energy finding ways not to.}

\end{enumerate}

\begin{baihoc}
Câu cà khịa chỉ có tác dụng khi người nghe nhận ra mình trong đó. Nếu em đọc và không thấy mình ở đâu cả --- đọc lại lần nữa.\\
\textit{(A roasting line only works when the listener recognizes themselves in it. If you read through and see yourself nowhere --- read it again.)}
\end{baihoc}
